\begin{textsample}{NEG Dim 1 – 2009\_c4 – Score -7.00 – 2009\_c4/enade\_2009\_biblio\_0e14e9ac55.txt}  \label{ex:f1_neg_030}
66 Capítulo 4 Percepção sobre a Prova As percepções dos estudantes sobre a prova aplicada no ENADE/2009 na área de BIBLIOTECONOMIA foram mensuradas por meio de nove itens que avaliaram desde o grau de dificuldade da prova até o tempo gasto para concluí -la.
As questões foram analisadas separando -se concluintes e ingressantes, e as percepções sobre a prova foram relacionadas com o desempenho dos estudantes e com a Região Geográfica de \textbf{localização} curso6.
O desempenho dos estudantes dentro de \textbf{cada} Área/Subárea foi classificado em quartos ( 1/4 ) para os concluintes e para os ingressantes em separado ( Ver figura 1 para o \textbf{esquema} de representação gráfica destas grandezas ).
As notas finais foram colocadas em \textbf{ordem} ascendente e no primeiro quarto foram colocadas as 25 \% notas mais baixas, limitadas superiormente pelo percentil 25 que vem a ser o primeiro quartil.
No segundo quarto, as 25 \% seguintes, limitadas superiormente pelo percentil 50 que vem a ser a mediana.
No terceiro quarto, as 25 \% notas após as encontradas no segundo quarto, ou seja, as primeiras 25 \% acima da mediana, limitadas superiormente pelo percentil 75 que vem a ser o terceiro quartil.
No último quarto foram alocadas as 25 \% maiores notas.
Potencialmente, poderia existir uma relação entre a nota do estudante e a sua percepção sobre a prova.
Calculando as médias das respostas segundo os quartos de notas, seria possível reconhecer a existência ( ou não ) de algum gradiente nas respostas, por exemplo, uma melhor nota correspondendo a uma percepção menos negativa.
Este gradiente seria traduzido, por exemplo, no gráfico 4.1a, pela diminuição das áreas laranja e vermelha ( Difícil e Muito difícil ) com o aumento do desempenho ( indo do primeiro ao último quarto ).
A medida tradicional de “ grau de discriminação ” comparou os resultados dos estudantes do primeiro quarto ( 25 \% piores estudantes ) aos do último quarto ( 25 \% melhores ).
Tal comparação, no gráfico disponibilizado, foi idêntica à comparação da primeira barra ( primeiro quarto ) com a última barra ( último quarto ) de \textbf{cada} grupo ingressante ou concluinte nos gráficos “ segundo quartos de nota final ” ( gráficos com números ímpares Por outro lado, a disponibilização das informações segundo os quartos de nota permitiu verificar o comportamento de dois outros grupos intermediários que na análise somente dos extremos tem que ser suposto monotônico.
A seguir, serão apresentados os principais resultados relativos aos nove itens avaliados. 4.1 Grau de dificuldade da prova 4.1.1 Formação Geral Os Gráficos 4.1a, 4.1b e 4.1c apresentam a distribuição das respostas à questão “ Qual o grau de dificuldade desta prova na parte de Formação Geral? ”, segundo quartos de nota final e grupo de estudantes concluintes e ingressantes que participaram do ENADE/2009 na área de BIBLIOTECONOMIA.
No total de concluintes, 16,0 \% declararam que a prova na parte da Formação Geral foi Difícil/Muito difícil, contra um valor semelhante para os ingressantes, 19,3 \%.
Tanto entre os concluintes quanto entre os ingressantes, a proporção de indivíduos que classificaram esta parte como Difícil/Muito difícil diminui com o desempenho do estudante, como era esperado.
Por outro lado, a proporção que declarou a parte Fácil/Muito fácil foi ligeiramente maior entre os concluintes ( 20,1 \% ) do que entre os ingressantes ( 19,0 \% ).
Não se observou uma variação monotônica significativa desta proporção com o aumento no desempenho, nem entre os concluintes nem entre os ingressantes.
Entre os concluintes, 63,9 \% classificaram esta parte com um nível médio de dificuldade.
A proporção entre os ingressantes foi ligeiramente menor, 61,7 \%.
Entre os concluintes, esta proporção aumentou com o desempenho, indo de 57,1 \% no primeiro quarto de notas a 67,0 \% no último quarto.
Os Gráficos 4.2a e 4.2b apresentam a distribuição das respostas à questão “ Qual o grau de dificuldade desta prova na parte de Formação Geral? ”, segundo a Região onde o curso \textbf{está} situado para os estudantes concluintes e ingressantes que participaram do ENADE/2009 na área de BIBLIOTECONOMIA.
Entre os concluintes, a maior proporção de estudantes que considerou a parte de Formação Geral da prova como de dificuldade média entre todas as Regiões brasileiras, foram os da Região Centro-Oeste com 67,0 \%, já entre os ingressantes, foram os estudantes da Região Nordeste com 69,1 \%.
A Região com as menores proporções desta resposta, foram respectivamente para os concluintes e ingressantes, a Sul e a Sudeste ( 59,3 \% e 55,1 \% ). 4.1.2 Componente Específico Os Gráficos 4.3a, 4.3b e 4.3c apresentam a distribuição das respostas à questão “ Qual o grau de dificuldade desta prova na parte do Componente Específico? ”, segundo quartos de nota final e grupo de estudantes concluintes e ingressantes que participaram do ENADE/2009 na área de BIBLIOTECONOMIA.
Em geral, para as desagregações consideradas, os estudantes classificaram o Componente Específico como sendo mais difícil que a Formação Geral.
No grupo de concluintes, 26,0 \% declararam que a prova na parte do Componente Específico foi Difícil/Muito difícil, sendo um valor muito maior para os ingressantes, 43,3 \%.
Como esperado, tanto entre concluintes quanto entre ingressantes, a proporção de indivíduos que classificaram esta parte como Muito Difícil diminui com o desempenho do estudante.
A proporção que declarou Difícil/Muito difícil entre os concluintes foi decrescente como função do desempenho, com exceção do segundo quarto, onde se verificou uma ascensão.
Entre os ingressantes, esta proporção foi decrescente e houve o desempenho com pico no segundo quarto.
A proporção que declarou a parte Fácil/Muito fácil foi maior entre os concluintes ( 10,1 \% ) do que entre os ingressantes ( 5,3 \% ).
Entre concluintes e ingressantes, a proporção dos que consideraram esta parte como Muito fácil decresceu com o desempenho.
Entre os concluintes, 63,9 \% classificaram esta parte com um nível médio de dificuldade.
A proporção entre os ingressantes foi menor, 51,45 \%.
Entre os concluintes, esta proporção basicamente aumentou com o desempenho ( com exceção do segundo quarto ), indo de 62,3 \% no primeiro quarto de notas a 67,0 \% no último quarto.
Entre os ingressantes, esta proporção foi crescente, indo de 47,6 \% no primeiro quarto a 60,5 \% no último.
Os Gráficos 4.4a e 4.4b apresentam a distribuição das respostas à questão “ Qual o grau de dificuldade desta prova na parte do Componente Específico? ”, segundo a Região onde o curso \textbf{está} situado para os estudantes concluintes e ingressantes que participam do ENADE/2009 na área de BIBLIOTECONOMIA.
Nenhum concluinte ou ingressante da Região Norte declarou que esta parte foi Muito fácil.
Entre os concluintes, a maior proporção de estudantes que considerou a parte do Componente Específico da prova como de dificuldade Os Gráficos 4.4a e 4.4b apresentam a distribuição das respostas à questão “ Qual o grau de dificuldade desta prova na parte do Componente Específico? ”, segundo a Região onde o curso \textbf{está} situado para os estudantes concluintes e ingressantes que participam do ENADE/2009 na área de BIBLIOTECONOMIA.
Nenhum concluinte ou ingressante da Região Norte declarou que esta parte foi Muito fácil.
Entre os concluintes, a maior proporção de estudantes que considerou a parte do Componente Específico da prova como de dificuldade média entre todas as Regiões brasileiras, foi na Região Norte com 71,83 \%, já entre os ingressantes, foram os estudantes da Região Nordeste com 56,7 \%.
A Região com as menores proporções desta resposta foram, respectivamente, para os concluintes e ingressantes, a Nordeste ( 61,2 \% ) e a Centro-Oeste ( 34,3 \% ). 4.2 Extensão da prova em relação ao tempo total Os Gráficos 4.5a, 4.5b e 4.5c apresentam a distribuição das respostas à questão sobre a extensão da prova em relação ao tempo total, segundo quartos de nota final e grupo de estudantes concluintes e ingressantes que participaram do ENADE/2009 na área de BIBLIOTECONOMIA.
Tanto para concluintes quanto para ingressantes, a maioria considerou a extensão da prova como adequada em relação ao tempo destinado à sua resolução, respectivamente 65,3 \% e 67,3 \%.
Entre concluintes e ingressantes esta proporção cresceu com o desempenho, com uma ligeira queda no último quarto.
A proporção dos que consideraram a extensão da prova Muito curta foi ligeiramente menor entre os concluintes ( 0,7 \% ) do que entre os ingressantes ( 0,8 \% ).
A proporção dos que consideraram a extensão da prova Curta/Muito curta vis-à-vis ao tempo total, foi menor entre os concluintes ( 3,3 \% ) do que entre os ingressantes ( 4,6 \% ).
A maior proporção de concluintes declarando que a prova foi Curta/Muito curta foi entre os estudantes de pior desempenho, 4,6 \%.
A proporção dos que consideraram a extensão da prova Muito longa foi próximo entre os concluintes ( 11,2 \% ) e ingressantes ( 10,8 \% ).
A proporção dos que consideraram a extensão da prova Longa foi ligeiramente menor entre os concluintes ( 21,9 \% ) do que entre os ingressantes ( 22,1 \% ).
Os Gráficos 4.6a e 4.6b apresentam a distribuição das respostas à questão sobre a extensão da prova em relação ao tempo total, segundo a Região onde o curso \textbf{está} situado para os estudantes concluintes e ingressantes que participam do ENADE/2009 na área de BIBLIOTECONOMIA.
Em todas as Regiões brasileiras, a maioria dos estudantes, tanto concluintes quanto ingressantes, perceberam a prova como adequada. 4.3 Compreensão dos Enunciados das Questões 4.3.1 Formação Geral Os Gráficos 4.7a, 4.7b e 4.7c apresentam a distribuição das respostas à questão “ Os enunciados das questões da prova na parte de Formação Geral, \textbf{estavam} claros e objetivos? ”, segundo quartos de nota final e grupo de estudantes concluintes e ingressantes que participaram do ENADE/2009 na área de BIBLIOTECONOMIA.
Com relação aos enunciados das questões da parte de Formação Geral, as opiniões foram positivas : a grande maioria dos concluintes ( 73,0 \% ) e ingressantes ( 73,9 \% ) consideraram claros e objetivos os enunciados de todas ou da maioria das questões.
Entre os concluintes e os ingressantes, esta proporção cresceu com o desempenho : 59,6 e 69,6 \% para o grupo com o pior desempenho e 78,3 \% e 81,9 \% para o de melhor, respectivamente.
Em linhas gerais, pode -se dizer que a combinação de “ Poucos ” e “ Não, nenhum ” decresceu com o desempenho tanto para concluintes como para ingressantes.
Os Gráficos 4.8a e 4.8b apresentam a distribuição das respostas à questão “ Os enunciados das questões da prova na parte de Formação Geral, \textbf{estavam} claros e objetivos? ”, segundo a Região onde o curso \textbf{está} situado para os estudantes concluintes e ingressantes que participaram do ENADE/2009 na área de BIBLIOTECONOMIA.
Em termos regionais, as proporções de ingressantes e de concluintes que consideraram claros e objetivos os enunciados de todas ou da maioria das questões não variaram muito.
Entre os concluintes, os valores foram de 69,2 \% no Sudeste a 78,8 \% no Nordeste.
Já entre os ingressantes, estas proporções variaram de 72,1 \% no Sudeste a 79,8 \% no Centro-Oeste.
Com relação às respostas “ Poucos ” e “ Não, nenhum ” tanto para concluintes quanto para ingressantes as maiores proporções ocorreram na Região Centro-Oeste ( 14,9 \% ) e Sudeste ( 13,7 \% ), respectivamente. 4.3.2 Componente Específico Os Gráficos 4.9a, 4.9b e 4.9c apresentam a distribuição das respostas à questão “ Os enunciados das questões da prova na parte de formação específica \textbf{estavam} claros e objetivos? ”, segundo quartos de nota final e grupo de estudantes concluintes e ingressantes que participaram do ENADE/2009 na área de BIBLIOTECONOMIA.
Com relação aos enunciados das questões da parte do Componente Específico, as opiniões foram também positivas : a grande maioria dos concluintes ( 71,1 \% ) e ingressantes ( 67,0 \% ) consideraram claros e objetivos os enunciados de todas ou da maioria das questões.
A proporção de concluintes e ingressantes declarando que todos os enunciados \textbf{estavam} claros e objetivos cresceu com o desempenho, sendo que entre os concluintes, há uma pequena queda no último quarto.
O número de estudantes que declararam que “ Não, nenhum ” dos enunciados das questões da parte do Componente Específico era claro e objetivo, foi equivalente com a parte de Formação Geral : 1,2 \% entre os concluintes e 1,5 \% entre os ingressantes.
Os que responderam “ Poucos ” são 10,5 \% entre os concluintes e 11,9 \% entre os ingressantes.
Os Gráficos 4.10a e 4.10b apresentam a distribuição das respostas à questão “ Os enunciados das questões da prova na parte de formação específica \textbf{estavam} claros e objetivos? ”, segundo a Região onde o curso \textbf{está} situado para os estudantes concluintes e ingressantes que participaram do ENADE/2009 na área de BIBLIOTECONOMIA.
Em termos regionais, as proporções de ingressantes e de concluintes que consideraram claros e objetivos os enunciados de todas ou da maioria das questões variaram entre os concluintes, de 67,8 \% no Sudeste a 77,1 \% no Nordeste e entre os ingressantes, de 63,5 \% no Nordeste a 76,9 no Centro-Oeste.
Com relação às respostas “ Poucos ” e “ Não, nenhum ” para ingressantes, a maior proporção ocorreu na Região Sudeste com 11,6 \%, enquanto que para os concluintes foi a Região Centro-Oeste com 9,0 \%. 4.4 Suficiência das informações/instruções fornecidas Os Gráficos 4.11a, 4.11b e 4.11c apresentam a distribuição das respostas à questão “ As informações/instruções fornecidas para a resolução das questões foram suficientes para resolvê -las? ”, segundo quartos de nota final e grupo de estudantes concluintes e ingressantes que participaram do ENADE/2009 na área de BIBLIOTECONOMIA.
Quando \textbf{perguntados} sobre as informações/instruções fornecidas para a resolução das questões, 81,1 \% dos concluintes e 77,5 \% dos ingressantes de todo o Brasil responderam que estas eram ou excessivas ou suficientes em todas ou na maioria das questões, demonstrando uma avaliação positiva da prova.
Levando -se em conta o desempenho dos participantes, percebeu -se consideráveis semelhanças entre as opiniões de estudantes, ingressantes e concluintes.
A proporção de estudantes com a opinião supracitada foi crescente com o desempenho, tanto para concluintes quanto para ingressantes.
Entre os concluintes, a proporção vai de 72,3 \% para os de pior desempenho a 85,3 \% para os de melhor desempenho.
Entre os ingressantes, as proporções equivalentes são, respectivamente 73,9 \% a 80,5 \%.
O complemento desta opinião, ou seja, os estudantes que responderam somente em algumas ou em nenhuma delas, apresentou uma proporção decrescente com o desempenho tanto para os concluintes quanto para os ingressantes.
A resposta “ Não, em nenhuma delas ” apareceu tanto para os concluintes quanto para os ingressantes em todos os grupos de desempenho.
Os Gráficos 4.12a e 4.12b apresentam a distribuição das respostas à questão “ As informações/instruções fornecidas para a resolução das questões foram suficientes para resolvê -las? ”, segundo a Região onde o curso \textbf{está} situado para os estudantes concluintes e ingressantes que participaram do ENADE/2009 na área de BIBLIOTECONOMIA. 4.5 Dificuldade encontrada ao responder à prova Os Gráficos 4.13a, 4.13b e 4.13c apresentam a distribuição das respostas à questão “ Você se deparou com alguma dificuldade ao responder a prova? ”, segundo quartos de nota final e grupo de estudantes concluintes e ingressantes que participaram do ENADE/2009 na área de BIBLIOTECONOMIA. \textbf{Perguntados} sobre as dificuldades com as quais se depararam ao responder à prova, 55,6 \% dos ingressantes apontaram como principal dificuldade o desconhecimento do conteúdo, já os 38,6 \% dos concluintes, apontaram como principal dificuldade a forma diferente de abordagem.
A forma diferente de abordagem do conteúdo foi considerada como maior dificuldade por 22,0 \% dos ingressantes, enquanto que a falta de motivação foi a segunda maior dificuldade para os concluintes, 27,1 \%, proporção que para os ingressantes é de 13,0 \%.
Esta dificuldade apresentou valores decrescentes conforme o desempenho para os concluintes.
Entre os ingressantes, esta dificuldade também apresenta valores decrescentes, com exceção do terceiro quarto que apresenta crescimento.
Uma proporção maior de concluintes do que de ingressantes declararam que não tiveram qualquer dificuldade para responder à prova, respectivamente 16,6 \% e 5,5 \%.
Os Gráficos 4.14a e 4.14b apresentam a distribuição das respostas à questão “ Você se deparou com alguma dificuldade ao responder a prova? ”, segundo a Região onde o curso \textbf{está} situado para os estudantes concluintes e ingressantes que participaram do ENADE/2009 na área de BIBLIOTECONOMIA.
Em relação às Regiões brasileiras, o percentual de ingressantes que apontaram o desconhecimento do conteúdo como principal Os Gráficos 4.14a e 4.14b apresentam a distribuição das respostas à questão “ Você se deparou com alguma dificuldade ao responder a prova? ”, segundo a Região onde o curso \textbf{está} situado para os estudantes concluintes e ingressantes que participaram do ENADE/2009 na área de BIBLIOTECONOMIA.
Em relação às Regiões brasileiras, o dificuldade ao responder à prova, variou de 47,6 \% na Região Sul, a 65,7 \% na Centro-Oeste.
Entre os concluintes, tal percentual variou de 6,4 \% na Região Centro-Oeste a 11,9 \%, na Sudeste.
A forma diferente de abordagem de conteúdo foi citada como maior dificuldade por 31,8 \% dos concluintes na Região Sudeste chegando, a 52,9 \% na Região Norte.
Já entre os ingressantes, os valores foram de 17,6 \% na Região Centro-Oeste a 26,0 \% na Região Nordeste.
O percentual de concluintes que citaram a falta de motivação como principal dificuldade variou de 16,52 \% na Região Norte a 32,3 \% na Região Sudeste.
Entre os ingressantes as estatísticas correspondentes foram de 6,3 \% na Região Norte a 16,5 \%, na Região Sudeste.
Os que declararam que não tiveram qualquer tipo de dificuldade variaram, entre os concluintes, de 11,8 \% na Região Norte a 20,65 \% na Região Sul.
Entre os ingressantes, este percentual variou de 3,7 \% na Região Centro-Oeste a 9,8 \% na Região Sul. 4.6 Percepção sobre os conteúdos das questões objetivas da prova Os Gráficos 4.15a, 4.15b e 4.15c apresentam a distribuição das respostas à questão “ Considerando apenas as questões objetivas da prova, você percebeu que :... ” segundo quartos de nota final e grupo de estudantes concluintes e ingressantes que participaram do ENADE/2009 na área de BIBLIOTECONOMIA.
As respostas possíveis para esta questão \textbf{estavam} relacionadas à \textbf{cobertura} do conteúdo das questões objetivas durante o curso.
Como esperado, a proporção de ingressantes ( 72,7 \% ) que declarou que não estudou ainda a maioria desses conteúdos foi bem maior do que entre os concluintes ( 4,1 \% ).
Esta proporção foi decrescente para os concluintes como função do desempenho.
No entanto, esta proporção para os ingressantes apresenta seus maiores percentuais no segundo e no terceiro grupo, desenhando um arco como função do desempenho.
Por outro lado, a proporção de estudantes que declarou que estudou e aprendeu todos ou a maioria desses conteúdos, entre os ingressantes foi de 13,1 \% e entre os concluintes 69,5 \%.
Esta proporção foi crescente como função do desempenho, tanto para os concluintes quanto para os ingressantes.
A proporção de estudantes que declarou ter estudado a maioria desses conteúdos, mas não os ter aprendido é bem maior entre os concluintes ( 17,6 \% ) do que os ingressantes ( 2,8 \% ).
Os Gráficos 4.16a e 4.16b apresentam a distribuição das respostas à questão “ Considerando apenas as questões objetivas da prova, você percebeu que :... ”, segundo a Região onde o curso \textbf{está} situado para os estudantes concluintes e ingressantes que participaram do ENADE/2009 na área de BIBLIOTECONOMIA.
Em relação às Regiões brasileiras, concluintes que afirmaram não ter estudado ainda a maioria dos conteúdos ocorreram em todas as Regiões, os percentuais variaram de 2,1 \% na Região Centro-Oeste a 5,0 \%, na Sudeste.
Entre os ingressantes, os percentuais variaram de 67,2 \% na Região Sul a 87,9 \%, na Centro-Oeste.
Em todas as Regiões, a maioria dos concluintes afirmou ter estudado e aprendido todos ou muito dos conteúdos, com percentuais que variaram de 61,4 \% na Região Norte a 80,9 \%, na Centro-Oeste.
Para os ingressantes, tal percentual variou de 5,6 \%, na Região Centro-Oeste, a 18,3 \%, na Sul. 4.7 Tempo gasto para concluir a prova Os Gráficos 4.17a, 4.17b e 4.17c apresentam a distribuição das respostas à questão “ Qual foi o tempo gasto para concluir a prova? ”, segundo quartos de nota final e grupo de estudantes concluintes e ingressantes que participaram do ENADE/2009 na área de BIBLIOTECONOMIA.
Quando \textbf{perguntados} quanto ao tempo de conclusão da prova, a maioria respondeu entre uma e duas horas, tanto entre os concluintes ( 40,1 \% ) quanto entre os ingressantes ( 41,4 \% ).
A segunda maior incidência de respostas foi entre duas e três horas, igualmente para ambos os grupos : concluintes ( 37,7 \% ) e ingressantes ( 35,8 \% ).
A proporção daqueles que declararam ter usado as quatro horas, mas não terem conseguido terminar foi bem baixo : 1,1 \% entre os ingressantes e 1,4 \% entre os concluintes.
Esta resposta ficou concentrada quase que exclusivamente entre os estudantes do primeiro quarto, com pior desempenho.
Entre os concluintes e os ingressantes percebeu -se uma tendência monotônica do tempo necessário para concluir a prova entre duas e três horas como função do desempenho.
Entre os ingressantes, a distribuição do tempo é em forma de U invertido para as respostas entre duas e três horas e entre três e quatro horas com valores extremos para estudantes com os melhores e piores desempenho.
Os Gráficos 4.18a e 4.18b apresentam a distribuição das respostas à questão “ Qual foi o tempo gasto para concluir a prova? ”, segundo a Região onde o curso \textbf{está} situado para os estudantes concluintes e ingressantes que participaram do ENADE/2009 na área de BIBLIOTECONOMIA.
A proporção de concluintes que gastaram entre uma e duas horas foi maior na Região Sudeste ( 44,5 \% ) e menor na Região Norte ( 19,0 \% ).
Para ingressantes, o maior percentual para este tempo ocorreu na Região Sudeste ( 44,3 \% ) e a menor na Região Centro-Oeste ( 29,4 \% ).
Entre os ingressantes, as Regiões correspondentes foram a Centro-Oeste com tempo mais longo e a Sul com tempo mais curto.

% matched lemmas: cada, cobertura, esquema, estar, localização, ordem, perguntar
\end{textsample}
